\subsection*{Updated objectives}

The primary objective is to implement and evaluate two unsupervised anomaly detection pipelines on MVTec 3D-AD under a unified experimental framework. The comparison focuses on predictive performance, localization quality, computational complexity and interpretability.

The updated research questions are:
\begin{itemize}
    \item How effective is a category-specific classical baseline based on raw normalized depth patches, PCA and One-Class SVM for 3D surface anomaly detection \cite{scholkopf2001estimating}?
    \item Does a compact deep learning model trained on normal samples learn more effective representations of surface regularity than local patch-based one-class methods \cite{li2025surveyindustrialad}?
    \item Does RGB information improve detection for defects that are weakly expressed in depth, such as contamination or color-related anomalies?
    \item What trade-offs emerge between predictive performance, localization quality, computational cost and interpretability when comparing classical, deep and multimodal variants?
\end{itemize}
